\icode{E-MK-SI-L1-260627-002}
在物理学的重大发现中，科学家们创造出了许多物理学研究方法，如理想实验法、极限思想法、类比法和建立物理模型法等等，以下关于所用物理学研究方法的叙述不正确的是\paren[A]
\begin{choices}
    \item 在不需要考虑物体本身的大小和形状时，用质点来代替物体的方法叫假设法
    \item 根据速度定义式\(v = \frac{\Delta x}{\Delta t}\)，当\(\Delta t\)非常非常小时，\(\frac{\Delta x}{\Delta t}\)就可以表示物体在\(t\)时刻的瞬时速度，该定义应用了极限思想法
    \item 在定义加速度\(a = \frac{\Delta v}{\Delta t}\)时，采用了比值定义法
    \item 在推导匀变速直线运动位移公式时，把整个运动过程划分成很多小段，每一小段近似看作匀速直线运动，然后把各小段的位移相加，这里采用了微元法
\end{choices}
\begin{answer}
    A\qquad\kaishu 将物体看作质点是理想模型法
\end{answer}